\documentclass[10pt,a4paper]{amsart}
\usepackage[margin=19mm]{geometry}
\usepackage{amsmath,amssymb,mathtools}
\usepackage[scr=rsfs,cal=euler]{mathalfa}
\usepackage{xfrac}
\usepackage[unicode,hidelinks,bookmarks=false]{hyperref}
\newcommand{\ie}{i.e.\ }
\newcommand{\cf}{cf.\ }
\newcommand{\resp}{respectively\ }
\newcommand{\Subsection}{\subsection}
\providecommand{\nobreakdash}{\nobreak\hbox{-}}
\setlength{\emergencystretch}{2em}
\multlinegap0pt
\usepackage{amssymb}
\usepackage{xcolor}
\usepackage{algorithm}\usepackage{algpseudocode}
\usepackage[normalem]{ulem}
\usepackage{tikz}
\newtheorem{proposition}{Proposition}
\def \be{{\mathbf {e}}}
\def \bx{{\mathbf {x}}}
\title{Optimal mean-time path planning for unmanned underwater vehicles: a Hamilton-Jacobi approach}
\author{Jonathan Valyou, Jeremy Brandman, Sanghyun Lee}
\date{Source: arXiv:2607.03407v1 (CC BY 4.0)}
\begin{document}
\maketitle

\noindent\emph{Source and licence.} Optimal mean-time path planning for unmanned underwater vehicles: a Hamilton-Jacobi approach. Version arXiv:2607.03407v1 (CC BY 4.0), distributed by arXiv under the Creative Commons Attribution 4.0 International licence, http://creativecommons.org/licenses/by/4.0/, as recorded in the arXiv licence metadata for this exact version and shown on the abstract page https://arxiv.org/abs/2607.03407v1. Full licence notice: \texttt{licenses/arxiv-2607.03407-CC-BY-4.0.txt}.\par\smallskip

\noindent\textit{Editorial scope.} Counted unit: the article's proposition eqn:unc-hj (source0.tex lines 453--467). The article's complete original proof of it is delivered with it (source0.tex lines 470--537, one \texttt{proof} environment of 68 lines). No mathematics is written, repaired or re-derived: the statement and the proof are the article's own lines, in the article's own notation, and no retained line is substituted or re-flowed.  No further source-local context is needed: every cross-reference of every retained line resolves inside the two delivered spans.  Nothing was omitted from any retained span, so the package carries no omission record.  No retained line contains a citation, so the wrapper carries no bibliography and no bibliography entry.  No retained line contains a figure environment, an image-inclusion command or drawing code, so no image is delivered and the package declares no asset dependency.  Editorial formatting only, in addition: an AMS \texttt{amsart} preamble that restates the article's own declarations and macros used by the retained lines, namely the environments \texttt{proposition} on separate counters, together with the standard packages those lines need.  The retained environments print in this document as Proposition 1 in order of appearance, whereas the article prints the same items with the numbering of its own full document.  The complete native source of the article as distributed in the arXiv e-print for this exact version is bundled unchanged as \texttt{sources/204383/source0.tex}.
\begin{proposition}
The minimum mean reachability time $u(\bx,\bx_S)$ satisfies the following system of Hamilton-Jacobi Partial Differential Equations:
%%%
\begin{align}
    & \min_{\theta \in  [0, 2\pi)} \{ -\nabla u \cdot s_{1,max} (\theta, \bx, T_1) \be_\theta + \sum_{i=1}^{n} p(i) \frac{s_{1,max} (\theta, \bx, T_1)}{s_{i,max} (\theta, \bx, T_i)} \} = 0 \label{eqn:unc-hj}\\
    & \nabla T_i(\bx, \bx_S,\gamma^*) \cdot s_{i,max}(\theta^*,\bx, T_i)\be_{\theta^*} = 1 , \quad \text{for } i=1,...,n
    \label{eqn:det-hj}
\end{align}
subject to the boundary conditions
\begin{align}
    & u(\bx_S, \bx_S) = 0 , \quad T_i(\bx_S, \bx_S, \gamma^*) = 0, \quad \text{for } i=1,...,n \label{eqn:bc-hj}
\end{align}
%%%
where the Hamiltonian minimizer is denoted $\theta^*$.
\end{proposition}

\begin{proof}

    Introduce a small change $dS$ in the optimal path connecting the starting point $x_S$ to another point $x$.  It follows from the dynamic programming principle and our definition of $u(\bx,\bx_S)$ that
    \begin{align*}
        &u(\bx,\bx_S) = \min\limits_{\substack{dS}}\{ u(\bx-dS,\bx_S) +  u(\bx,\bx-dS)\}.
    \end{align*}
    

    
    Define $0<t_1<t_2<...<t_{p-1}<t_p<...<T_f$ with uniform time-step size $\Delta t=t_p-t_{p-1}$.  Let $dS=\Delta t \cdot s_1(\theta,\bx,T_1)\be_\theta$ where $\Delta t \ll 1$; it follows that
    
    \begin{align*}
        &u(\bx,\bx_S) = \min\limits_{\substack{\theta \in [0,2\pi) \\ 0 {\color{violet} <} s_1(\theta,\bx,T_1) \leq s_{1,max}(\theta, \bx,T_1)}} \{ u(\bx-\Delta t \cdot s_1(\theta,\bx,T_1)\be_\theta, \bx_S) + u(\bx,\bx-\Delta t \cdot s_1(\theta,\bx,T_1)\be_\theta) \}.
    \end{align*}
    By applying a Taylor Expansion to the first term on the right-hand side, we obtain
    \begin{equation}
        u(\bx,\bx_S) = \min\limits_{\substack{\theta \in [0,2\pi) \\ 0 {\color{violet} <}  s_1(\theta,\bx,T_1) \leq s_{1,max}(\theta, \bx,T_1)}} \{ u(\bx,\bx_{S})- \nabla u(\bx,\bx_{S}) \cdot \Delta t \cdot s_1(\theta,\bx,T_1)\be_\theta + u(\bx,\bx-\Delta t \cdot s_1(\theta,\bx,T_1)\be_\theta)\}.
    \end{equation}
    After canceling $u(\bx,\bx_S)$ from each side, we get
    \begin{equation}
        0 = \min\limits_{\substack{\theta \in [0,2\pi) \\ 0 {\color{violet} <}  s_1(\theta,\bx,T_1) \leq s_{1,max}(\theta, \bx,T_1)}} \{ - \nabla u \cdot \Delta t \cdot s_1(\theta,\bx,T_1) \be_\theta + u(\bx,\bx-\Delta t \cdot s_1(\theta,\bx,T_1)\be_\theta)\}.
    \end{equation}


    Next, since $\Delta t\ll1$, $\Delta t \cdot s_1(\theta,\bx,T_1)\be_\theta$ is small.  Therefore, it follows that the optimal path from $\bx-\Delta t \cdot s_1(\theta,\bx,T_1)\be_\theta$ to $\bx$ is approximately straight in the direction $\be_\theta$ with a length of $\Delta t \cdot s_1(\theta,\bx,T_1)$.  Additionally, note that the velocity under the $i$-th ocean model would be $s_i(\theta,\bx,T_i)\be_\theta$.  We recover the travel time under the $i$-th ocean model by taking the path length divided by the $i$-th model velocity resulting in $\Delta t \frac{s_1(\theta,\bx,T_1)}{s_i(\theta,\bx,T_i)}$.  Recall, that we are trying to find the mean optimal time $u$ along this path so we must apply the probability distribution of realizing each velocity under the path resulting in the reachability time expression

    \begin{equation}
        u(\bx,\bx-\Delta t \cdot s_1(\theta,\bx,T_1)\be_\theta) = \Delta t \cdot \sum_{i=1}^{n} p(i) \frac{s_1 (\theta, \bx, T_1)}{s_i (\theta, \bx, T_i)}),
    \end{equation}
    which substituting into the above equation results in 
    \begin{equation}
        0 = \min\limits_{\substack{\theta \in [0,2\pi) \\ 0 {\color{violet} <}  s_i(\theta,\bx,T_i) \leq s_{i,max}(\theta, \bx,T_i) }} \{ - \nabla u \cdot \Delta t \cdot s_1(\theta,\bx,T_1)\be_\theta + \Delta t \cdot \sum_{i=1}^{n} p(i) \frac{s_1 (\theta, \bx, T_1)}{s_i (\theta, \bx, T_i)})\}.
    \end{equation}
    %%
    Observe that the right-hand side of the above expression is linear in $s_1(\theta,\bx,T_1)$ so it achieves its minimum at the endpoint or the largest possible value of $s_1(\theta,\bx,T_1)$.  Since we assume an optimal path exists, it follows that $s_1(\theta,\bx,T_1)=s_{1,max}(\theta,\bx,T_1)$ minimizes the right side of the equation.  Additionally, we note that the right-hand side of the above expression varies inversely with respect to each $s_i(\theta,\bx,T_i)$.  It follows that setting $s_i(\theta,\bx,T_i)=s_{i,max}(\theta,\bx,T_i)$ minimizes the right side of the equation and yields
    \begin{equation}
        0 = \min\limits_{\substack{\theta \in [0,2\pi) }} \{ - \nabla u \cdot \Delta t \cdot s_{1,max}(\theta,\bx,T_1)\be_\theta + \Delta t \cdot \sum_{i=1}^{n} p(i) \frac{s_{1,max} (\theta, \bx, T_1)}{s_{i,max} (\theta, \bx, T_i)})\}.
    \end{equation}
To complete the derivation, we divide by $\Delta t$ and obtain the governing PDE
 \eqref{eqn:unc-hj}:
    \begin{align}
        0 = \min\limits_{\substack{\theta \in [0,2\pi)}} \{ - \nabla u  \cdot s_{1,max}(\theta,\bx,T_1)\be_\theta + \sum_{i=1}^{n} p(i) \frac{s_{1,max} (\theta, \bx, T_1)}{s_{i,max} (\theta, \bx, T_i)}\}.
    \end{align}
%%
Lastly, note that the above argument indicates that the optimal control direction $\theta^*$ corresponds to the Hamiltonian minimizer $\theta$ in \eqref{eqn:unc-hj}.  From the control directions $\theta^*$, the optimal path $\gamma^*$ from $\bx_S$ to an arbitrary target position $\bx$ is obtained.

In order to derive the other PDEs in \eqref{eqn:det-hj}, we take a similar approach. As before, we use the Dynamic Programming Principle to break the path, now $\gamma^*$, into two parts.  However, in this case we use the choice $dS = \Delta t \cdot s_{1,max}(\theta^*,\bx,T_1)\be_{\theta^*}$
made in the derivation above, to arrive at  
\begin{equation}
        T_i(\bx,\bx_S,\gamma^*) \approx T_i(\bx-\Delta t \cdot (s_{i,max}(\theta^*,\bx,T_i)\be_{\theta^*}),\bx_S,\gamma^*) + \Delta t.
\end{equation}
By applying a Taylor Expansion to the right-hand side, we get
\begin{equation}
        T_i(\bx,\bx_S,\gamma^*) = T_i(\bx,\bx_S,\gamma^*) - \Delta t \cdot \nabla T_i(\bx,\bx_S,\gamma^*) \cdot (s_{i,max}(\theta^*,\bx,T_i)\be_{\theta^*}) + \Delta t.
\end{equation}
Subtracting $T_i(\bx,\bx_S,\gamma^*)$ from each side results in
\begin{equation}
        0 = - \Delta t \cdot \nabla T_i(\bx,\bx_S,\gamma^*) \cdot (s_{i,max}(\theta^*,\bx,T_i)\be_{\theta^*}) + \Delta t.
\end{equation}
Again, division by $\Delta t$ results in
    \begin{equation}
        0 =- \nabla T_i(\bx,\bx_S,\gamma^*) \cdot (s_{i,max}(\theta^*,\bx,T_i)\be_{\theta^*}) + 1.
    \end{equation}
Finally, we can rewrite the above equation to yield \eqref{eqn:det-hj}:
    \begin{equation}
        1 = \nabla T_i(\bx,\bx_S,\gamma^*) \cdot (s_{i,max}(\theta^*,\bx,T_i)\be_{\theta^*}).
    \end{equation}
\end{proof}

\end{document}
